\chapter{Conclusiones}
\label{ch:conclusiones}

El presente documento ha presentado la \textbf{Revisión Crítica de Diseño (CDR)}
del subsistema de \textbf{Super--Resolución basada en Inteligencia Artificial
(SAI)} del proyecto CubeSat INTISAT, concebido como un módulo complementario de
postprocesamiento para las imágenes adquiridas por el \textit{CMOS Microscopy
Payload (CMOSM)}.

A lo largo del documento se ha definido un alcance técnico claro y conservador,
alineado con principios de ingeniería de sistemas ECSS/NASA, evitando atribuir
al subsistema capacidades que excedan su función real dentro de la arquitectura
de la misión.

\section{Síntesis de resultados}

El análisis realizado permite concluir que:

\begin{itemize}
    \item El subsistema SAI cumple su rol como herramienta de apoyo visual para la
    interpretación de imágenes microscópicas adquiridas en configuración de
    microscopía sin lentes.
    \item La arquitectura del pipeline de super--resolución es modular, trazable
    y coherente con la separación entre adquisición física y procesamiento
    computacional.
    \item Los modelos de super--resolución evaluados permiten mejorar la
    legibilidad perceptual de estructuras ya presentes en la imagen original,
    sin alterar la información base adquirida.
    \item La estrategia de verificación aplicada demuestra el correcto
    funcionamiento del subsistema dentro de un marco experimental y
    cualitativo.
\end{itemize}

\section{Integración con el CMOS Microscopy Payload}

El subsistema SAI se integra de manera natural con el CMOSM, manteniendo una
separación estricta entre:

\begin{itemize}
    \item adquisición de datos primarios (CMOSM),
    \item procesamiento y realce visual posterior (SAI).
\end{itemize}

Esta separación garantiza que la validez científica del sistema de adquisición
no depende del uso de técnicas de inteligencia artificial, preservando la
trazabilidad y la reproducibilidad de los resultados experimentales.

\section{Limitaciones reconocidas}

Se reconoce explícitamente que:

\begin{itemize}
    \item el subsistema SAI no incrementa la resolución óptica real del sistema,
    \item los resultados generados no son adecuados para análisis cuantitativos
    ni metrológicos,
    \item la inferencia se realiza fuera de línea y no constituye una función
    crítica de misión.
\end{itemize}

Estas limitaciones han sido incorporadas de manera consciente en el diseño y la
documentación del subsistema.

\section{Estado de madurez del subsistema}

En base a la evidencia presentada, el subsistema SAI alcanza un nivel de madurez
adecuado para su integración como \textbf{demostrador tecnológico}, permitiendo:

\begin{itemize}
    \item la evaluación controlada de técnicas modernas de super--resolución,
    \item la generación de productos derivados para apoyo visual,
    \item la exploración de futuras extensiones metodológicas.
\end{itemize}

\section{Trabajo futuro}

Como líneas de trabajo futuras se identifican:

\begin{itemize}
    \item ampliación del conjunto de datos de evaluación,
    \item comparación sistemática entre distintos modelos de super--resolución,
    \item estudio del impacto de la super--resolución en distintos tipos de
    muestras,
    \item evaluación preliminar de viabilidad de ejecución en plataformas
    embebidas, sin comprometer la arquitectura actual.
\end{itemize}

\section{Conclusión final}

La Revisión Crítica de Diseño del subsistema SAI concluye que el diseño propuesto
es técnicamente consistente, correctamente acotado y compatible con los
objetivos del proyecto INTISAT. El enfoque adoptado permite incorporar técnicas
de inteligencia artificial de forma responsable, transparente y defendible,
fortaleciendo el valor experimental y formativo de la misión sin comprometer la
integridad científica del sistema de adquisición.
